\section{Graph limits, permutons, and the Schr\"odinger problem}\label{app:limits}

This appendix explains in more detail the analytic objects used in
\cref{sec:limits,sec:heuristic,sec:deterministic}. Its aim is to let a reader who has not
met graphons or entropic transport check each step. Statements marked as theorems are
quoted from the literature or proved here. Heuristic steps are labelled as such.

\subsection{Dense graph limits}
For a graph $G$ on $[n]$, define the step function $W_G$ on $[0,1]^2$ to be $1$ on
$\bigl(\frac{i-1}n,\frac in\bigr]\times\bigl(\frac{j-1}n,\frac jn\bigr]$ when $i\sim j$ and
$0$ otherwise. Thus $W_G$ is ``the adjacency matrix drawn as a picture''. A
\emph{graphon} is a symmetric measurable $W:[0,1]^2\to[0,1]$. For a fixed graph $F$
on $[k]$ the homomorphism density is
\[
t(F,W)=\int_{[0,1]^k}\prod_{\{u,v\}\in E(F)}W(x_u,x_v)\,dx_1\cdots dx_k .
\]
$t(F,W_G)$ is the probability that a random map $V(F)\to V(G)$ is a homomorphism. The
\emph{cut norm} is $\|W\|_\square=\sup_{S,T}\bigl|\int_{S\times T}W\bigr|$, the supremum
over measurable $S,T\subseteq[0,1]$. The \emph{cut distance}
$\delta_\square(W,U)=\inf_\phi\|W-U^\phi\|_\square$ takes the infimum over
measure-preserving relabellings $\phi$ of $[0,1]$.

\begin{theorem}[Lov\'asz--Szegedy \cite{LovaszSzegedy}; see \cite{Lovasz}]
Every sequence of graphs with $|V(G_n)|\to\infty$ has a subsequence along which
$t(F,W_{G_n})$ converges for every $F$. The limits are $t(F,W)$ for a graphon $W$, and
convergence of all $t(F,\cdot)$ is equivalent to $\delta_\square(W_{G_n},W)\to0$.
\end{theorem}

Conversely, a graphon $W$ generates \emph{$W$-random graphs}: choose
$x_1,\dots,x_n$ independently and uniformly from $[0,1]$, and join $i,j$ independently with
probability $W(x_i,x_j)$. These converge to $W$ almost surely. For example,
$G(n,\tfrac12)\to\tfrac12$, and the graphs with $i\sim j\iff i+j\le n$ converge to
$\mathbf 1[x+y\le1]$.

\subsection{Sparse graphs, and why $G(n)$ needs more}
If $e(G_n)=o(n^2)$ then $W_{G_n}\to0$, which carries no information. The theory of $L^p$
graphons \cite{BCCZ1,BCCZ2} rescales by the edge density $\rho_n=2e(G_n)/n^2$. One
studies $W_{G_n}/\rho_n$, whose limits are symmetric $W\in L^p([0,1]^2)$, $p>1$, with no
upper bound required. Compactness requires \emph{$L^p$ upper regularity}: block averages of
$W_{G_n}/\rho_n$ over partitions into parts of measure at least $\eta$ must have $L^p$ norm
at most a constant $K$ independent of $\eta$. The sampling model joins $i,j$ with
probability $\min(1,\rho_nW(x_i,x_j))$.

For $G(n)$ we have $\rho_n\asymp n^{-1/2}$. The natural candidate limit is $\kappa(x,y)=1/(2\sqrt{x+y})$:
the squares near $m$ have density $1/(2\sqrt m)$, and
$\sqrt n\cdot 1/(2\sqrt{n(x+y)})=\kappa(x,y)$. \Cref{prop:mod4} shows that cut-norm
convergence to $\kappa$ fails, and the remark after it shows that upper regularity fails.
Both failures come from residues: whether $i+j$ is a square depends on $i\bmod q$ and
$j\bmod q$ for every $q$. The limit must therefore record residues.

\subsection{The profinite integers}
$\Zhat=\varprojlim\mathbb{Z}/q\mathbb{Z}\cong\prod_p\mathbb{Z}_p$ is a compact ring
containing $\mathbb{Z}$ densely. An element is a compatible choice of residues modulo every
$q$. Its Haar probability measure $\mu_{\Zhat}$ gives each residue class modulo $q$ mass
$1/q$. Let $\sigma$ be the law of $k^2$ when $k$ is Haar-distributed. With
$\varrho_q(a)=\#\{x\bmod q:\ x^2\equiv a\}$,
\[
\sigma(\{s\equiv a\bmod q\})=\varrho_q(a)/q .
\]
The squares in $\mathbb{Z}_p$ have Haar measure $\sum_{k\ge0}p^{-2k}\cdot\frac{p-1}{2p}=\frac{p}{2(p+1)}$
for odd $p$, and $\sum_{k\ge0}4^{-k}\cdot\frac18=\frac16$ for $p=2$. The infinite
product over $p$ is $0$, so the squares form a Haar-null set in $\Zhat$, and $\sigma$
is singular with respect to $\mu_{\Zhat}$.

\begin{proposition}[joint equidistribution]\label{prop:equid}
Fix $q$, residues $r,s\bmod q$, and intervals $I,J\subseteq[\delta,1]$, $\delta>0$. Then
\[
\#\bigl\{(i,j):\ i\sim j,\ \tfrac in\in I,\ \tfrac jn\in J,\ i\equiv r,\ j\equiv s\ (\mathrm{mod}\ q)\bigr\}
=\frac{\varrho_q(r+s)}{q^2}\,n^{3/2}\!\iint_{I\times J}\kappa+O_{q,\delta}(n).
\]
\end{proposition}

\begin{proof}
Fix $i\equiv r$. Its neighbours are $j=k^2-i$, and $j\equiv s$ if and only if
$k^2\equiv r+s\pmod q$. Within any $q$ consecutive values of $k$ there are exactly
$\varrho_q(r+s)$ such $k$. Apply \cref{lem:squaresum} with $f=\mathbf 1_J$ to the
subsequence of $k$ in a fixed residue class; this changes the mesh by the factor $q$.
The count for fixed $i$ is $\frac{\varrho_q(r+s)}q\sqrt n\int_J\kappa(i/n,y)\,dy+O_{q,\delta}(1)$.
Now sum over the $\frac nq|I|+O(1)$ values of $i$.
\end{proof}

Accordingly, put $\Omega=[0,1]\times\Zhat$ with measure $\lambda\otimes\mu_{\Zhat}$, and send
vertex $i$ to $(i/n,i)$. Define the operator
\[
(Tf)(x,\alpha)=\int_0^1\kappa(x,y)\,\mathbb{E}_{s\sim\sigma}\,f(y,\,s-\alpha)\,dy .
\]
\Cref{prop:equid} says exactly that the normalised adjacency operator
$(A_nf)(i)=n^{-1/2}\sum_{j\sim i}f(j)$ converges to $T$ against all test functions that
depend on $\alpha$ only through a fixed modulus. In the language of Backhausz--Szegedy
\cite{BackhauszSzegedy}, a \emph{graphop} is a positivity-preserving self-adjoint
operator $L^\infty(\Omega)\to L^1(\Omega)$. \emph{Action convergence} asks that, for
every $k$ and all $f_1,\dots,f_k$ with $\|f_i\|_\infty\le1$, the joint laws of
the $2k$-tuples
\[(f_1,\dots,f_k,\,A_nf_1,\dots,A_nf_k)\]
be approximated by those for $T$, and conversely.
\Cref{conj:limit} asserts this for $T$. None of the arguments below depends on it. Its
role is to explain why a single analytic object, $\kappa$ together with $\sigma$,
should govern all of the asymptotics.

\subsection{Permutons and their large deviations}
A \emph{permuton} is a probability measure on $[0,1]^2$ with both marginals uniform; the
set of permutons is $\Pi(\lambda,\lambda)$. A permutation $\pi\in S_n$ defines the
permuton $\mu_\pi$ with density $n$ on the squares
$\bigl(\frac{i-1}n,\frac in\bigr]\times\bigl(\frac{\pi(i)-1}n,\frac{\pi(i)}n\bigr]$. We use the weak
topology on permutons.

\begin{theorem}[Trashorras; Kenyon--Kr\'al'--Radin--Winkler {\cite[Thm.~1]{KKRW}}]\label{thm:LDPperm}
If $\pi$ is uniform in $S_n$, the laws of $\mu_\pi$ satisfy a large deviation principle
with speed $n$ and rate function $I(\mu)=\KL(\mu\,\|\,\lambda^2)=\int\rho\log\rho$ if $\mu$
has density $\rho$, and $I=+\infty$ otherwise. Informally,
$\#\{\pi:\mu_\pi\approx\mu\}\approx n!\,e^{-n\KL(\mu\|\lambda^2)}$.
\end{theorem}

By Varadhan's lemma \cite[Thm.~4.3.1]{DemboZeitouni}, for bounded continuous $\Phi$ on
permutons,
\[
\lim_{n\to\infty}\frac1n\log\mathbb{E}_\pi\,e^{n\Phi(\mu_\pi)}=\sup_\mu\{\Phi(\mu)-I(\mu)\}.
\]
We apply this to $\Phi(\mu)=\int\log\kappa\,d\mu$. Since $\kappa\ge 1/(2\sqrt2)$,
$\log\kappa$ is bounded below. It is unbounded above only at the corner, where
$\kappa(x,y)\le 1/(2\sqrt x)$. Let $\kappa_M=\min(\kappa,M)$. Only rows with
$i<n/(4M^2)$ can have $\kappa>M$, so for every permutation
\[
0\le\sum_i\Bigl[\log\kappa\bigl(\tfrac in,\tfrac{\pi(i)}n\bigr)-\log\kappa_M\bigl(\tfrac in,\tfrac{\pi(i)}n\bigr)\Bigr]
\le\sum_{i<n/(4M^2)}\tfrac12\log\frac{n}{4M^2i}=O\!\left(\frac n{M^2}\right).
\]
Also $\sum_i\log\kappa_M(i/n,\pi(i)/n)=n\int\log\kappa_M\,d\mu_\pi+o(n)$ uniformly in $\pi$,
by uniform continuity. Letting $M\to\infty$ after $n\to\infty$ proves:

\begin{proposition}\label{prop:firstmoment}
Let $R_n$ be the random $0/1$ matrix with independent entries
$\mathbb{P}(R_n(i,j)=1)=p_{ij}=\frac1{2\sqrt{i+j}}$. Then
\[
\mathbb{E}\operatorname{per}R_n=\sum_{\pi\in S_n}\prod_ip_{i\pi(i)}=\frac{n!}{n^{n/2}}\ \mathbb{E}_\pi\prod_i\kappa\bigl(\tfrac in,\tfrac{\pi(i)}n\bigr),
\qquad
\frac1n\log\mathbb{E}_\pi\prod_i\kappa\bigl(\tfrac in,\tfrac{\pi(i)}n\bigr)\longrightarrow C,
\]
where $C=\sup_{\mu\in\Pi(\lambda,\lambda)}\{\int\log\kappa\,d\mu-\KL(\mu\|\lambda^2)\}$.
\end{proposition}

This is a theorem, given \cref{thm:LDPperm}. The \emph{heuristic} steps of
\cref{sec:heuristic} are: (i) the passage from $\mathbb{E}\operatorname{per}R_n$, the
first moment of the number of cycle covers of the random model, to Hamiltonian cycles,
i.e.\ to cyclic $\pi$, at a cost of $e^{o(n)}$; and (ii) the passage from $R_n$ to $G(n)$.
\Cref{sec:deterministic} replaces (ii) for cycle covers by a deterministic argument.

\subsection{The Schr\"odinger problem}
We now analyse the variational problem
\[
C=\sup_{\mu\in\Pi(\lambda,\lambda)}\Bigl\{\int\log\kappa\,d\mu-\KL(\mu\,\|\,\lambda^2)\Bigr\}.
\]

\emph{Gibbs variational principle.} For a probability measure $\nu$ and measurable $g$
with $\int e^g\,d\nu<\infty$,
\[
\sup_\mu\Bigl\{\int g\,d\mu-\KL(\mu\|\nu)\Bigr\}=\log\int e^g\,d\nu,
\]
attained exactly at $d\mu\propto e^g\,d\nu$. This follows from Jensen's inequality applied to
$\log\int e^{g}\frac{d\nu}{d\mu}\,d\mu$. Without the marginal constraints, our
supremum would therefore be $\log\iint\kappa=\log\frac43(\sqrt2-1)$. The constraints are what
make the problem a \emph{Schr\"odinger bridge}: with
$R=\kappa\,dx\,dy/\iint\kappa$,
\[
C=\log\iint\kappa-\min_{\mu\in\Pi(\lambda,\lambda)}\KL(\mu\,\|\,R).
\]
This is the static Schr\"odinger problem, equivalently entropic optimal transport with cost
$c=-\log\kappa$ and regularisation parameter $1$ \cite{Leonard,Nutz,PeyreCuturi}.

\emph{Existence and uniqueness.} $\Pi(\lambda,\lambda)$ is weakly compact, $\KL(\cdot\|R)$
is lower semicontinuous and strictly convex, and $\KL(\mu\|R)<\infty$ for $\mu=\lambda^2$.
So there is a unique minimiser $\mu^*$.

\emph{Form of the minimiser.} Perturb $\mu$ within $\Pi(\lambda,\lambda)$ and introduce
Lagrange multipliers $\varphi(x),\psi(y)$ for the two marginal constraints. Stationarity of
$\int\log\kappa\,d\mu-\int\rho\log\rho$ gives $\log\kappa-\log\rho-1=\varphi(x)+\psi(y)$,
that is,
\[
d\mu^*=a(x)\,b(y)\,\kappa(x,y)\,dx\,dy .
\]
The rigorous version, the product form of $I$-projections onto sets of measures with
given marginals, is due to Csisz\'ar \cite{Csiszar}. Borwein--Lewis--Nussbaum
\cite{BLN} treat kernels in exactly this ``$DAD$'' form, including measurable kernels. Since
$\kappa$ is symmetric and $\mu^*$ is unique, $\mu^*$ is symmetric. Then $a(x)b(y)=a(y)b(x)$,
so $b$ is a constant multiple of $a$, and after rescaling $b=a$. The marginal constraints
become the \emph{Schr\"odinger equation}
\begin{equation}\label{eq:schrod}
a(x)\int_0^1a(y)\,\kappa(x,y)\,dy=1\qquad(0\le x\le1).
\end{equation}

\emph{The value.} Since $\log\frac{d\mu^*}{d\lambda^2}=\log a(x)+\log a(y)+\log\kappa$, we get
$\int\log\kappa\,d\mu^*-\KL(\mu^*\|\lambda^2)=-\int(\log a(x)+\log a(y))\,d\mu^*=-2\int_0^1\log a$,
using the uniform marginals. Hence $C=-2\int_0^1\log a$.

\emph{A dual formula giving certified bounds.} For any measurable $u>0$ and any
$\mu\in\Pi(\lambda,\lambda)$, apply the Gibbs principle with $g=\log(\kappa\,u\otimes u)$, and use
$\int\log(u\otimes u)\,d\mu=2\int\log u$:
\[
\int\log\kappa\,d\mu-\KL(\mu\|\lambda^2)\ \le\ \log\iint\kappa(x,y)u(x)u(y)\,dx\,dy-2\int_0^1\log u .
\]
Equality holds for $u=a$ by \eqref{eq:schrod}. So
\[
C=\min_{u>0}\Bigl\{\log\iint\kappa\,u\otimes u-2\int\log u\Bigr\}
 =\max_{\mu\in\Pi(\lambda,\lambda)}\Bigl\{\int\log\kappa\,d\mu-\KL(\mu\|\lambda^2)\Bigr\}.
\]
Any trial function $u$ gives an upper bound for $C$, and any permuton gives a lower bound.
With $u=1$ we get $C\le\log\iint\kappa=-0.5937$. The comparison $C<B$ of
\cref{prop:bregman} is the same principle with only the first marginal constrained.

\subsection{Sinkhorn's iteration}
Write $(Kf)(x)=\int_0^1\kappa(x,y)f(y)\,dy$. Sinkhorn's iteration
\cite{Sinkhorn,SinkhornKnopp} alternately enforces the two marginals:
$b_{k+1}=1/(Ka_k)$, $a_{k+1}=1/(Kb_{k+1})$. A fixed point satisfies $a\,Kb=1=b\,Ka$. In the
symmetric case we use the damped form
\[
a_{k+1}=\bigl(a_k/Ka_k\bigr)^{1/2},
\]
the geometric mean of $a_k$ and the Sinkhorn update. Its fixed points are exactly the
solutions of \eqref{eq:schrod}.

\emph{Convergence.} On positive functions, use Hilbert's projective metric
$d_H(f,g)=\log\bigl(\sup\frac fg\cdot\sup\frac gf\bigr)$. The map $f\mapsto1/f$ is a
$d_H$-isometry. By the Birkhoff--Hopf theorem \cite{Birkhoff}, a positive linear operator
whose image has projective diameter $\Delta<\infty$ contracts $d_H$ by the factor
$\tanh(\Delta/4)$. Taking geometric means halves distances. So the damped iteration
contracts with factor at most $\frac12\bigl(1+\tanh(\Delta/4)\bigr)<1$ whenever
$\kappa$ is bounded above and below, and then $\Delta\le2\log(\sup\kappa/\inf\kappa)$.
Our $\kappa$ is unbounded only at the origin. For the truncation $\kappa_M$ the iteration
converges geometrically. $C$ is monotone in the kernel, and the estimate preceding
\cref{prop:firstmoment} gives $0\le C-C_M=O(M^{-2})$. Nussbaum \cite{Nussbaum} treats
the scaling of integral kernels in this Hilbert-metric framework in much greater generality.

\emph{Numerics.} With the midpoint rule on $M$ nodes, $K$ becomes an $M\times M$ matrix.
For $M=400,1600,3200$ the iteration converges to residual below $10^{-12}$ and gives
$C=-0.632547$, $-0.632539$ and $-0.632538$ respectively (\cref{alg:sinkhorn-cont}). The potential
$a$ increases from about $0.82$ at $x=0$, a finite value because $\kappa(0,\cdot)$ is
integrable, to about $1.76$ at $x=1$. It is plotted in the talk figures.

\subsection{The discrete counterpart: scaling, permanents, Bethe}\label{app:bethe}
The finite-$n$ analogue of a permuton is a doubly stochastic matrix. The analogue of the
Schr\"odinger problem is matrix scaling.

\emph{Matrix scaling.} If a nonnegative matrix $A$ has total support, there are positive
diagonal $D_1,D_2$ with $D_1AD_2$ doubly stochastic (Sinkhorn--Knopp
\cite{SinkhornKnopp}), and for symmetric $A$ one may take $D_1=D_2=D=\mathrm{diag}(d_i)$.
Then $\operatorname{per}A=\operatorname{per}(S)/\prod_id_i^2$ with $S=DAD$, and the van der
Waerden bound $\operatorname{per}S\ge n!/n^n$ \cite{Egorychev,Falikman} gives
\[
\log\operatorname{per}A\ \ge\ \log\frac{n!}{n^n}-2\sum_i\log d_i=-n-2\sum_i\log d_i+O(\log n).
\]

\emph{Bethe.} \Cref{thm:GS} with $P=S$ gives
$\log\operatorname{per}A\ge F(S)=-2\sum_i\log d_i-n+O\bigl(\sum s_{ij}^2\bigr)$. Indeed
$-\sum s_{ij}\log s_{ij}=-\sum s_{ij}(\log d_i+\log d_j)=-2\sum_i\log d_i$ because the rows and
columns sum to $1$, and $\sum(1-s)\log(1-s)=-\sum s+O(\sum s^2)=-n+O(\sum s^2)$. For
$G(n)$, $\sum s_{ij}^2\le n\max s_{ij}=O(\sqrt n)$. So van der Waerden and Bethe-at-$S$
agree up to $o(n)$. The advantage of \cref{thm:GS} is that \emph{any} doubly stochastic $P$
is allowed, so an exact scaling is not needed. This is what makes \cref{prob:L1} a
perturbation problem.

\emph{Discrete versus continuum.} Suppose $d_i=n^{-1/4}a(i/n)(1+o(1))$ in an averaged sense.
Then
\[
-2\sum_i\log d_i=\frac n2\log n-2\sum_i\log a(i/n)+o(n)=\frac n2\log n+Cn+o(n),
\]
and hence $F(S)=\frac n2\log n+(C-1)n+o(n)$. This is exactly the first-moment value of
\cref{prop:firstmoment}, since $\log n!-\frac n2\log n=\frac n2\log n-n+O(\log n)$. The
ansatz $n^{1/4}d_i\approx a(i/n)$ is consistent because the scaling equation
$d_i\sum_{j\sim i}d_j=1$ becomes, by \cref{lem:squaresum},
$a(x)\cdot\sqrt n\int a\,\kappa(x,\cdot)\cdot n^{-1/2}=1$, which is \eqref{eq:schrod}. The
same computation for the explicit ansatz \eqref{eq:ansatz} gives
$-\sum\widetilde p\log\widetilde p=\frac12\log n\sum\widetilde p-2\sum_i\log a(i/n)\sum_j\widetilde p_{ij}
=\frac n2\log n+Cn+o(n)$, since its row sums are $1+o(1)$ on average.

\emph{Where the residues went.} In the limit on $\Omega$, the analogue of
\eqref{eq:schrod} uses the kernel $\kappa(x,y)$ times the arithmetic coupling $\nu$ of
$(\alpha,\beta)$, $\beta=s-\alpha$. If $\alpha$ is Haar-distributed and $s\sim\sigma$ is
independent, then $\beta$ is Haar-distributed. So $\nu$ \emph{already has uniform
marginals}, the potentials do not depend on $\alpha$, and the optimal coupling on
$\Omega^2$ is $\mu^*\otimes\nu$. The arithmetic layer contributes $0$ to $C$. In finite
terms, this is why \cref{lem:squaresum} needs no congruence conditions: the row sums of
the ansatz involve only the sizes of the neighbours. The numerical values of $F(S)$ in
\cref{sec:deterministic}, computed for the actual graphs $G(n)$, confirm this.

\subsection{Dictionary}
\begin{center}\small
\begin{tabular}{ll}
\toprule
continuum (limit object) & finite $n$ (the graph $G(n)$)\\
\midrule
kernel $\kappa(x,y)$ on $[0,1]^2$ & $\sqrt n\,A_{G(n)}$, via \cref{lem:squaresum}\\
arithmetic coupling $\nu$ on $\Zhat^2$ & residues of $i,j$ with $i+j$ a square, \cref{prop:equid}\\
permuton $\mu\in\Pi(\lambda,\lambda)$ & doubly stochastic $P$ supported on $E(G(n))$\\
$\int\log\kappa\,d\mu-\KL(\mu\|\lambda^2)$ & $F(P)-\frac n2\log n+n$, the Bethe free energy\\
Schr\"odinger potential $a(x)$ & $n^{1/4}d_i$, the scaling vector\\
Sinkhorn on $L^\infty[0,1]$ & Sinkhorn on $A_{G(n)}$\\
$C$ & $\frac1n\bigl(\log\operatorname{per}A_{G(n)}-\frac n2\log n\bigr)+1$ (conjecturally)\\
$B=\int\log f$ (Gibbs, one marginal) & Bregman's bound via the degrees $d(i)$\\
\bottomrule
\end{tabular}
\end{center}
